\documentclass[a4paper,12pt]{article}
\usepackage{graphicx}
\usepackage{float} 
\usepackage{wrapfig}
\usepackage{amsmath}
\usepackage[table]{xcolor}
\usepackage[hidelinks]{hyperref}
\graphicspath{{./img/}}

\title{Définition des règles et Architecture du projet}
\author{Anissa KOURBAN HOUSSEN FASSY \\ Yohann FRONT-REIGNIER \\ Hajar CHERQI \\ He YU}
\date{30 avril 2025}

\begin{document}
\maketitle
\renewcommand{\contentsname}{Sommaire}
\tableofcontents


\section{Définition des règles}
\subsection{Règles générales}
Le jeu se joue à 2 joueurs, noir et blanc, sur une grille de 8x8 (de 1 à 8 et de A à H) composée de carrés. Il se déroule au tour par tour et débute avec 2 pions de chaque couleur comme illustré dans la Figure 1 : 2 pions noirs sont placés en E4 et D5, tandis que 2 pions blancs sont placés en D4 et E5.
\begin{figure}[ht]
    \centering
    \includegraphics[width=0.25\textwidth]{start_game.png}
    \caption{Position initiale des pions}
\end{figure}

\subsection{Déroulé d'une partie}
Noir commence toujours la partie en premier. Pour capturer des pions adverses, il faut les entourer avec des pions à nous. Si un joueur ne peut pas capturer de pions, il passe son tour. Les pions capturés sont tous ceux situés entre les 2 premiers pions qui se "croisent" comme illustré dans la Figure 2.
\begin{figure}[ht]
    \centering
    \includegraphics[width=0.25\textwidth]{first_choice.png}
    \caption{Premier choix de capture du joueur noir}
\end{figure}

\subsection{Conditions de fin de partie}
\begin{itemize}
    \item Aucun joueur ne peut jouer.
    \item L'othellier est plein.
\end{itemize}
\begin{figure}[ht]
    \centering
    \includegraphics[width=0.26\textwidth]{last_choice.png}
    \caption{Exemple de fin de partie}
\end{figure}

Le choix du vainqueur est fait en fonction du joueur ayant le plus de pions à lui sur l'othellier.

\section{Architecture du projet}
Notre projet est réalisé en Python car ce langage est, pour nous, le plus adapté pour faire un jeu comme Othello. Il permet une programmation orientée objet, ce qui facilite la gestion des règles, des pions, et de l'interface graphique.

\subsection{Classes principales}
\subsubsection{Classe \emph{Othellier}}
La classe \emph{Othellier} représente la grille du jeu et gère les règles principales. Parmi les méthodes qu’elle contient, on retrouve notamment:
\begin{itemize}
    \item \textbf{placePawn(position)} : Place un pion sur la grille si la position est valide.
    \item \textbf{calculatePossibilities()} : Calcule tous les coups possibles pour le joueur actuel.
    \item \textbf{winner()} : Détermine le gagnant en fonction du nombre de pions.
    \item \textbf{\_findOpponentPawns(x, y)} : Trouve les pions adverses situés entre le pion en entrée et les autres pions du joueur.
    \item \textbf{\_takePawns(pawn1, pawn2)} : Capture les pions adverses entre deux positions.
    \item \textbf{\_checkDirection(x,y,dx,dx)} : Vérifie les diagonales autour d'un pion de coordonnées x,y.
    \item \textbf{notPlayable()} : Vérifie si aucun des deux joueurs ne peuvent jouer.
    \item \textbf {game\_over()}:
    vérifie si nous avons atteint la fin de la partie.
\end{itemize}

\subsubsection{Classe \emph{Pawn}}
La classe \emph{Pawn} représente un pion sur l'othellier. Elle est définie par :
\begin{itemize}
    \item \textbf{color} : La couleur du pion (0 pour noir, 1 pour blanc).
    \item \textbf{position} : La position du pion sur la grille.
\end{itemize}
Elle contient les méthodes suivantes :
\begin{itemize}
    \item \textbf{getColor()} : Retourne la couleur du pion.
    \item \textbf{changeColor()} : Change la couleur du pion.
    \item \textbf{getPosition()} : Retourne la position du pion.
\end{itemize}

\subsubsection{Classe \emph{Computer}}
La classe \emph{Computer} hérite de \emph{Othellier} et implémente des stratégies pour jouer contre un humain. Elle utilise l'algorithme Minimax avec élagage Alpha-Beta pour choisir le meilleur coup. Elle contient les méthodes suivantes :
\begin{itemize}
    \item \textbf{evaluation(othellier)} : Prend en entrée une instance de la classe Othellier et renvoie une valeur ( notion expliquer dans 2.3.2).
    \item \textbf{minimax(othellier,profondeur,alpha,beta)}:
    appliquer les principes de l'algorithme minimax (notion expliquer en 3.2.1) et renvoie le pion le plus avantageux à jouer.
\end{itemize}

\subsubsection{Interface graphique}
L'interface graphique est gérée par les classes suivantes :
\begin{itemize}
    \item \textbf{Interface} : La fenêtre principale du jeu.
    \item \textbf{MainMenu} : Le menu principal avec les options de jeu.
    \item \textbf{Game} : La classe qui gère l'affichage de la grille et des pions pendant la partie.
    \item \textbf{GameOver} : L'écran de fin de partie.
\end{itemize}

\subsection{Représentation de l'othellier}
L'othellier est représenté par une matrice 8x8 contenant des objets \emph{Pawn}. Voici un exemple de matrice initiale :
\[
\begin{bmatrix}
    0 & 0 & 0 & 0 & 0 & 0 & 0 & 0\\
    0 & 0 & 0 & 0 & 0 & 0 & 0 & 0\\
    0 & 0 & 0 & 0 & 0 & 0 & 0 & 0\\
    0 & 0 & 0 & \text{Pawn}(0, (3,3)) & \text{Pawn}(1, (3,4)) & 0 & 0 & 0\\
    0 & 0 & 0 & \text{Pawn}(1, (4,3)) & \text{Pawn}(0, (4,4)) & 0 & 0 & 0\\
    0 & 0 & 0 & 0 & 0 & 0 & 0 & 0\\
    0 & 0 & 0 & 0 & 0 & 0 & 0 & 0\\
    0 & 0 & 0 & 0 & 0 & 0 & 0 & 0
\end{bmatrix}
\]

\subsection{Joueur contre Ordinateur}
Pour la partie en Joueur contre Ordinateur, nous avons implémenté l'algorithme Minimax avec élagage Alpha-Beta. Cet algorithme est utilisé pour simuler plusieurs coups à l'avance et choisir le meilleur coup possible pour l'ordinateur. Voici une explication détaillée de son fonctionnement :

\subsubsection{Principe de l'algorithme Minimax}
L'algorithme Minimax est une méthode récursive utilisée pour prendre des décisions optimales dans des jeux à deux joueurs. Il repose sur les principes suivants :
\begin{itemize}
    \item \textbf{Maximisation} : Le joueur (ordinateur) cherche à maximiser son score.
    \item \textbf{Minimisation} : L'adversaire cherche à minimiser le score de l'ordinateur.
    \item \textbf{Évaluation des coups} : Chaque coup est évalué en fonction d'une fonction de score qui prend en compte des critères comme le nombre de pions capturés, la position sur la grille, etc.


En entrée, la fonction prend deux paramètres : l’othellier (état actuel du plateau) et une profondeur. Cette profondeur est décrémentée à chaque itération récursive de l’algorithme.

Le déroulement de la fonction est le suivant :

Vérification des conditions d’arrêt
À chaque appel, on commence par vérifier si :
la partie est terminée (plus aucun coup possible pour les deux joueurs),
ou si la profondeur est égale à 1.

   Dans notre implémentation, le joueur 0 représente l’humain, tandis que le joueur 1 représente la machine. Étant donné que l’objectif de l’algorithme est de trouver le meilleur coup pour la machine, la profondeur initiale doit être impair pour que la dernière évaluation se fasse au tour de la machine.

Application de la simulation
Si les conditions d’arrêt ne sont pas remplies, on passe à la simulation :
On initialise une variable meilleur score à -infini si le joueur actuel est la machine (maximisation), ou à +infini si c’est l’humain (minimisation).
On vérifie que le joueur actuel dispose de coups possibles.
Pour chaque coup possible :
On crée une copie de l’othellier pour ne pas modifier l’état réel du jeu.
On place le pion du joueur à la position simulée.
On appelle récursivement la fonction Minimax avec le nouvel othellier, l’adversaire comme joueur, et une profondeur décrémentée.

Mise à jour du meilleur coup
À chaque retour de la fonction récursive, on compare la valeur retournée avec la meilleure valeur trouvée jusqu’à présent :
Si la nouvelle valeur est meilleure (selon que l’on maximise ou minimise), on la conserve.

Retour du résultat
Une fois tous les coups possibles explorés, la fonction retourne le meilleur coup à jouer pour la machine, c’est-à-dire celui qui lui donne l’avantage stratégique maximal.

\subsubsection{Fonction d'évaluation}
La fonction d’évaluation utilisée dans notre jeu s’inspire des stratégies qu’un joueur humain appliquerait lorsqu’il doit choisir la meilleure position pour placer son pion. Ce choix doit être réfléchi, car plusieurs éléments influencent la qualité d’un coup :
\begin{itemize}

 \item \textbf{La différence de pions  entre les deux joueurs } : un indicateur simple mais utile, notamment en fin de partie.
 \item \textbf{La valeur stratégique des cases} : chaque case de l’othellier a une importance différente. Pour cela, nous utilisons une matrice 8x8 dans laquelle chaque case est associée à un poids représentant sa valeur stratégique (par exemple, les coins valent beaucoup plus que les cases centrales). Notre fonction évalue la somme des poids des cases occupées par le joueur.
\item \textbf{Le contrôle des coins} : les coins (comme A1, H1, A8, H8) sont essentiels, car une fois acquis, ils ne peuvent plus être repris par l’adversaire. De plus, ils permettent de sécuriser les pions voisins et d’en capturer plus facilement.
\end{itemize}
Notre fonction d’évaluation combine ces trois critères de la manière suivante :
La différence entre le nombre de pions du joueur et de l’adversaire ;
La somme des poids des cases occupées par le joueur selon la grille stratégique ;
Le nombre de coins contrôlés par le joueur.

Chacune de ces composantes est multipliée par un coefficient pondérateur, déterminé selon son importance stratégique. La somme des coefficients est égale à 1. La fonction retourne alors la somme pondérée de ces trois éléments, fournissant ainsi une estimation de la qualité d’un coup.

\end{itemize}



\subsubsection{Élagage Alpha-Beta}
L'élagage Alpha-Beta est une optimisation de l'algorithme Minimax qui permet de réduire le nombre de coups à explorer. Il fonctionne en éliminant les branches de l'arbre de recherche qui ne peuvent pas influencer la décision finale. Cela améliore considérablement les performances, surtout pour des jeux comme Othello où le nombre de coups possibles peut être élevé.

\begin{itemize}
    \item \textbf{Alpha} : La meilleure valeur que le joueur maximisant (ordinateur) peut garantir à ce stade.
    \item \textbf{Beta} : La meilleure valeur que le joueur minimisant (adversaire) peut garantir à ce stade.
    \item Si à un moment donné, Alpha devient supérieur ou égal à Beta, la branche correspondante est ignorée (coupure).
\end{itemize}

Lorsqu’un joueur maximise (la machine), la valeur alpha est mise à jour comme suit :
alpha = max(alpha, score)
Cela signifie que si l’on découvre un score supérieur à la meilleure valeur actuelle, on le retient.
Si ce score dépasse ou atteint la valeur de beta, cela signifie que le joueur minimisant (adversaire) n’aurait jamais laissé cette situation se produire. On peut donc couper cette branche, car elle ne sera jamais choisie : c’est une coupure alpha-beta.

Lorsqu’un joueur minimise (l’humain), la valeur beta est mise à jour ainsi :
beta = min(beta, score)

Cela signifie que si un score plus faible est trouvé, on l’enregistre comme nouvelle borne.
De la même manière, si beta devient inférieur ou égal à alpha, cela indique que le joueur maximisant (la machine) ne retiendra jamais ce chemin. On peut donc ignorer cette branche également.

Grâce à cette technique d’élagage, l'algorithme évite d'explorer des branches inutiles de l’arbre de recherche, ce qui permet :
une réduction significative de la complexité du Minimax ,
une exécution plus rapide, en particulier dans un jeu comme Othello où le nombre de coups possibles peut être élevé à chaque tour.

\subsubsection{Application dans Othello}
Dans notre implémentation, l'algorithme Minimax avec élagage Alpha-Beta est utilisé pour :
\begin{itemize}
    \item \textbf{Simuler plusieurs coups à l'avance} : L'ordinateur évalue les conséquences de ses coups et anticipe les réponses possibles de l'adversaire.
    \item \textbf{Valoriser les coins de la grille} : Les coins sont stratégiquement importants car ils ne peuvent pas être capturés une fois occupés.
    \item \textbf{Pénaliser les cases adjacentes aux coins} : Ces cases sont dangereuses car elles permettent à l'adversaire de capturer un coin.
    \item \textbf{Évaluer les positions} : Une fonction d'évaluation calcule un score basé sur le nombre de pions capturés, les positions stratégiques, et le contrôle de la grille.
\end{itemize}

\subsubsection{Exemple de fonctionnement}
Supposons que l'ordinateur joue en tant que Noir et que la grille est dans l'état suivant :
\[
\begin{bmatrix}
    0 & 0 & 0 & 0 & 0 & 0 & 0 & 0\\
    0 & 0 & 0 & 0 & 0 & 0 & 0 & 0\\
    0 & 0 & 0 & 0 & 0 & 0 & 0 & 0\\
    0 & 0 & 0 & \text{Pawn}(0, (3,3)) & \text{Pawn}(1, (3,4)) & 0 & 0 & 0\\
    0 & 0 & 0 & \text{Pawn}(1, (4,3)) & \text{Pawn}(0, (4,4)) & 0 & 0 & 0\\
    0 & 0 & 0 & 0 & 0 & 0 & 0 & 0\\
    0 & 0 & 0 & 0 & 0 & 0 & 0 & 0\\
    0 & 0 & 0 & 0 & 0 & 0 & 0 & 0
\end{bmatrix}
\]

L'algorithme Minimax avec élagage Alpha-Beta fonctionne comme suit :
\begin{enumerate}
    \item L'ordinateur génère tous les coups possibles pour Noir.
    \item Pour chaque coup, il simule les réponses possibles de l'adversaire (Blanc).
    \item Il évalue chaque scénario en utilisant une fonction de score.
    \item Il choisit le coup qui maximise son score tout en minimisant les opportunités de l'adversaire.
\end{enumerate}

\subsubsection{Performances}
Grâce à l'élagage Alpha-Beta, l'algorithme peut explorer efficacement les coups possibles sans examiner toutes les branches de l'arbre de recherche. Cela permet à l'ordinateur de jouer rapidement même lorsque la grille est presque pleine.

\begin{table}[H]
    \centering
    \renewcommand{\arraystretch}{2} % Augmente l'espacement entre les lignes
    \begin{tabular}{|c|c|}
        \hline
        \rowcolor{blue!40} \textbf{Profondeur} & \textbf{Temps moyen de décision} \\
        \hline
        \rowcolor{blue!30} 3 & 0.80s \\
        \hline
        \rowcolor{blue!30} 5 & 1s 10 \\
        \hline
        \rowcolor{blue!30} 7 & 1s 50 \\
        \hline
        \rowcolor{blue!30} 9 & 3s 30 \\
        \hline
    \end{tabular}
\caption{Temps moyen de prise de décision par l'IA selon la profondeur}
\end{table}


D’après le tableau ci-dessus, on en déduit que la profondeur idéale se situe autour de 5 ou 7.
Ces valeurs permettent de traiter un nombre suffisant de cas tout en conservant un temps de calcul raisonnable.
Une profondeur de 3 est certes très rapide, mais elle n’explore pas suffisamment de branches, ce qui limite l'efficacité de l'IA.
À l’inverse, une profondeur de 9 offre une meilleure anticipation, mais le temps de calcul devient trop élevé, malgré l’utilisation de l’élagage Alpha-Beta.


\subsubsection{Limites}
Bien que l'algorithme soit performant, il a certaines limites :
\begin{itemize}
    \item \textbf{Profondeur de recherche limitée} : L'algorithme ne peut pas explorer toutes les possibilités si la profondeur de recherche est trop grande. Lorsque la profondeur donnée est trop grande l'algorithme ralenti et met trop de temps à jouer son coup.
    \item \textbf{Fonction d'évaluation} : La qualité des décisions dépend fortement de la précision de la fonction d'évaluation.
\end{itemize}
\subsection{Interface Graphique}

Sur l'écran d'accueil quatre boutons sont disponibles :
\begin{itemize}
    \item \textbf{Joueur vs Joueur} : lance le jeu en mode joueur contre joueur.
    \item \textbf{Joueur vs IA} : lance le jeu en mode joueur contre intelligence artificielle.
    \item \textbf{Règles du jeu} : affiche les règles du jeu ainsi que les options proposées pour aider le joueur (comme des indications sur les cases jouables).
    \item \textbf{Quitter} : ferme la fenêtre.
\end{itemize}

Lorsque l’on clique sur les boutons \textit{Joueur vs Joueur} ou \textit{Joueur vs IA}, l’interface de jeu apparaît. On y trouve :
\begin{itemize}
    \item un plateau (l’othellier),
    \item une indication du nombre de points de chaque joueur affichée de chaque côté,
    \item une indication du joueur dont c’est le tour,
    \item deux boutons :
    \begin{itemize}
        \item un pour afficher les possibilités de jeu pour le tour en cours,
        \item un autre pour les afficher tout au long de la partie (un nouveau clic permet de désactiver cette option).
    \end{itemize}
\end{itemize}

À la fin de la partie, l’interface de fin de jeu apparaît, annonce le vainqueur, et propose deux boutons :
\begin{itemize}
    \item \textbf{Rejouer} : renvoie vers l’écran d’accueil,
    \item \textbf{Quitter} : ferme la fenêtre.
\end{itemize}


\section{Diagramme UML}
Le diagramme UML suivant illustre les relations entre les différentes classes du projet :
\begin{figure}[ht!]
    \centering
    \includegraphics[scale=0.5]{OthelloDiagram.png}
    \caption{Diagramme du projet Othello}
\end{figure}

\section{Test}

Nous avons développé un module de test permettant de vérifier le bon fonctionnement des fonctions essentielles de notre code, afin de garantir la stabilité et la cohérence du programme.

Les tests portent notamment sur les classes \texttt{Othellier} et \texttt{Pawn} :
\begin{itemize}
    \item Vérification du changement de joueur : il est essentiel que le changement de joueur s’effectue correctement. Une erreur à ce niveau compromettrait les calculs et le déroulement des parties.
    \item Vérification du placement des pions : le programme doit s’assurer qu’un pion est bien posé lorsqu’un joueur en fait la demande. Si le tour change sans que le pion ait été placé, cela endommage la logique de la partie.
    \item Vérification de la fonction de calcul des coups possibles : cette fonction est cruciale pour déterminer les actions autorisées à chaque tour.
    \item Vérification d’autres fonctions critiques : nous testons également diverses fonctions dont le dysfonctionnement pourrait compromettre le déroulement normal de la partie.
\end{itemize}

\section{Difficultés rencontrés}

Au cours de la construction du projet, nous avons rencontré plusieurs problèmes qui ont été résolus par la suite. Ces difficultés nous ont permis d’aborder les obstacles sous un autre angle et de renforcer notre travail en équipe.

\begin{itemize}
    \item \textbf{Lien entre l'interface et les modules \texttt{Othellier} et \texttt{Pawn}} :  
    La communication entre l’interface graphique et les modules principaux s’est révélée parfois complexe. Cela nous a amenés à revoir notre implémentation et à adapter notre vision du code afin d'améliorer la structure globale du programme.

    \item \textbf{Implémentation de l’algorithme Minimax} :  
    L'intégration de l'algorithme Minimax a soulevé plusieurs difficultés :
    \begin{itemize}
        \item la gestion de la récursion, 
        \item la simulation de coups sans altérer l’état réel de la matrice de l’\texttt{Othellier}, ce qui nous a obligés à manipuler des copies de la matrice pour préserver l'état du plateau,
        \item la gestion de la complexité algorithmique et du temps d'exécution.
    \end{itemize}
    Pour surmonter ces défis, nous nous sommes documentés sur le fonctionnement de Minimax et ses variantes, ainsi que sur les stratégies employées par les meilleurs joueurs. Cela nous a permis de développer une intelligence artificielle capable de proposer un réel défi et d’offrir une expérience de jeu agréable.
\end{itemize}


\end{document}